\section{INTRODUCTION}
\label{sec:intro}

In practical deployment scenarios, due to acquisition costs, equipment setup constraints, and limited camera trajectories, it is often difficult to obtain multi-view images with sufficient coverage. Compared with dense-view capture, sparse inputs provide limited information about visible regions, causing 3D reconstruction methods to overfit local observations and struggle to maintain global geometric consistency~\cite{Yang23_FreeNeRF}. Sparse observations also introduce greater geometric ambiguity, making occluded regions, textureless areas, and view boundaries prone to missing geometry, structural distortions, and degraded appearance reconstruction quality~\cite{Niemeyer22_RegNeRF}. Establishing stable geometric initialization and completing unseen regions remain critical challenges in sparse-view 3D reconstruction.

To improve reconstruction efficiency, Kerbl et al.~\cite{Kerbl23_3DGS} introduced 3D Gaussian Splatting (3DGS), an explicit scene representation and rendering framework. 3DGS represents scene geometry and appearance with differentiable Gaussian ellipsoids, whose parameters are jointly optimized. With efficient tiled rasterization, 3DGS enables rapid training and real-time rendering~\cite{Kerbl23_3DGS}, providing advantages over implicit neural representations. Subsequent variants further improve structured scene representation and alias-free rendering~\cite{Lu24_ScaffoldGS,Yu24_MipSplatting}.

However, under sparse-view conditions, Structure-from-Motion (SfM)~\cite{Schonberger16_COLMAP} often recovers sparse and incomplete point clouds. Since 3DGS typically relies on SfM-calibrated camera parameters and sparse point clouds for initialization, insufficient initialization affects optimization stability and geometric convergence of Gaussian primitives~\cite{Chung24_DepthRegGS}. The lack of cross-view constraints can further cause geometric degradation, floaters, and pseudo-geometry~\cite{Xu25_DropoutGS}. During optimization, sparse observations provide limited supervision~\cite{Park25_DropGaussian}, and visibility imbalance further aggravates overfitting~\cite{Somraj23_ViPNeRF}, degrading novel-view rendering quality.

Building upon 3DGS, existing sparse-view methods have explored depth regularization~\cite{Li24_DNGaussian} and densification~\cite{Zhu24_FSGS}. Others strengthen sparse-view training via co-regularization or Gaussian dropout~\cite{Zhang24_CoRGS,Park25_DropGaussian}. Most 3DGS-based sparse reconstruction methods rely on depth supervision and predefined sampling strategies for Gaussian initialization. Although these methods improve reconstruction quality, they still face challenges in capturing detailed structures and maintaining geometric consistency under sparse observations.

To address this issue, we present SynGS, a 3D Gaussian Splatting method guided by geometric and depth priors for sparse-view 3D reconstruction. SynGS strengthens Gaussian initialization through explicit geometric constraints and depth-guided densification, while Dynamic Visibility Regularization reduces overfitting during optimization. The main contributions are summarized as follows:
\begin{itemize}
\item \textbf{Visual Hull Initialization.} We construct a visual hull from multi-view silhouettes to generate a stable geometric prior for Gaussian initialization.
\item \textbf{VGGT / Depth-Guided Densification.} We use cross-view consistent depth predictions to back-project supplementary points, supplementing missing geometry and improving Gaussian initialization.
\item \textbf{Dynamic Visibility Regularization.} We introduce dynamic visibility regularization during training by randomly masking Gaussian primitives to improve gradient distribution and alleviate sparse-view overfitting.
\end{itemize}
